\documentclass[11pt]{article}
\usepackage[T1]{fontenc}
\usepackage{lmodern}
\usepackage[margin=1in]{geometry}
\usepackage{amsmath,amssymb,amsthm,mathtools}
\usepackage{microtype}
\usepackage{enumitem}
\usepackage{booktabs}
\usepackage[hidelinks]{hyperref}

\newtheorem{theorem}{Theorem}
\newtheorem{lemma}[theorem]{Lemma}
\newtheorem{proposition}[theorem]{Proposition}
\newtheorem{corollary}[theorem]{Corollary}
\theoremstyle{definition}
\newtheorem{definition}{Definition}
\newtheorem{example}{Example}
\theoremstyle{remark}
\newtheorem*{remark}{Remark}

\newcommand{\N}{\mathbb N}
\newcommand{\Q}{\mathbb Q}
\newcommand{\V}{\mathcal V}
\newcommand{\W}{\mathcal W}
\newcommand{\Hist}{\mathcal H}
\newcommand{\Dist}{\Delta_{\Q}}
\newcommand{\enc}{\operatorname{enc}}
\newcommand{\dom}{\operatorname{dom}}
\newcommand{\sgn}{\operatorname{sgn}}
\newcommand{\id}{\mathrm{id}}
\newcommand{\tagword}{\operatorname{tag}}
\newcommand{\Prob}{\mathbb P}
\newcommand{\Exp}{\mathbb E}
\newcommand{\one}{\mathbf 1}
\newcommand{\eps}{\varepsilon}
\newcommand{\TV}{\mathrm{TV}}
\newcommand{\paren}[1]{\left(#1\right)}
\newcommand{\norm}[1]{\left\lVert#1\right\rVert}

\title{A symmetric universal machine with\texorpdfstring{\\}{ }identically zero universal intelligence}
\author{}
\date{17 November 2026}

\begin{document}
\maketitle

\begin{abstract}
Alexander and Hutter ask whether a reward--punishment symmetric universal
machine can make the sign of universal intelligence computable.  For every
fixed suffix-free reinforcement-learning encoding, we construct an ordinary
prefix-free universal Turing machine for which the original intelligence
sum is identically zero.  Symmetry holds for every computable environment,
and the zero identity holds for every policy, whether computable or not.
The environment class is unchanged: only the limits of finite expected
reward totals must exist in the prescribed bounded interval for every
policy.  Neither absolute convergence of rewards nor a uniform bound on
intermediate totals is assumed.

The construction has two parts.  A computable injective transformation of
full environment kernels reverses every policy's value and reflects
convergence.  A zero-value probe preserves even the entries at histories
that the original environment cannot reach.  Finite inverse tests and a
contracting triangular structure leave one zero-valued kernel in the
infinite backward core.  Two partial program branches then implement the
two parities of finite backward depth.  Equal-length tags pair every
complete description, including competing universal descriptions, and
therefore pair actual shortest-description weights.  The resulting
constant-zero sign algorithm answers the stated existence question.  A
requirement that the intelligence functional distinguish policies would
be an additional condition.
\end{abstract}

\section{Introduction}
\label{sec:introduction}

Universal intelligence evaluates a policy by its performance across a
class of computable environments, with simpler environments receiving
larger weights.  In the formulation of Legg and Hutter
\cite{LeggHutter2007}, these weights are determined by shortest programs
in a fixed universal language.  Alexander and Hutter
\cite{AlexanderHutter2021} extend the setting to rewards of either sign.
An environment may reward or punish a policy, and the limiting expected
total reward may be positive, zero, or negative.

This extension introduces a natural symmetry.  Reverse the signs of all
rewards in both the history and the next output of an environment.  A
universal language is symmetric when this operation preserves the
shortest description length of every computable environment.  Such
languages make the intelligence score change sign when a policy's
interpretation of rewards is reversed.  The question in
\cite[Section~5]{AlexanderHutter2021} is whether a symmetric universal
language can make the sign of the score computable from a computable
policy.  The strong version asks for the sign at zero as well as at
nonzero values; the weak version permits arbitrary behavior when the
score is zero.

The answer proved here is affirmative for the strong version.  For any
fixed suffix-free encoding of histories and rational probability
distributions, there is a symmetric ordinary universal machine whose
intelligence score is zero for every policy.  The same machine still
simulates every prefix-free partial computable binary function.  Every
well-behaved computable environment still appears in the score with a
strictly positive shortest-description weight.  The zero result comes
from an exact pairing of their contributions.

The conclusion concerns the existence question as stated.  An identically
zero score does not rank policies or measure differences in their
performance.  It shows that symmetry and universality alone do not
exclude this outcome.  Asking for a nonzero or policy-discriminating
score with computable sign imposes a further condition, which the present
construction does not settle.

\subsection*{Why cancellation needs a construction}

Ordinary reward reversal is insufficient.  If $D$ denotes reward duality,
the basic identity is
\[
  V_{D\mu}^{\pi}=-V_{\mu}^{D\pi}.
\]
The policy on the right has changed.  A policy may react to its observed
rewards, so this identity does not pair opposite values for one fixed
policy.  Cancellation for every fixed policy requires a different
operation on environments.

The limiting-value convention also matters.  An allowed environment
need only have convergent \emph{finite expected reward totals}.  Its
intermediate totals may make arbitrarily large excursions, and the sum
of absolute rewards need not be integrable.  A construction that assumes
an absolute tail bound would address a smaller class.  We instead use
a random cutoff, refund the finite visible prefix, and make a long
negative payment based on one privately sampled reward.  A specific
geometric identity proves that the transformed totals converge exactly
when the original totals do.

The sum is over full kernels.  Two environments count as different when
they disagree at a history that neither can reach from the initial
state.  A transformation based only on reachable behavior can lose this
distinction, preventing a pairing of the original summands.  Our
construction includes a separate experiment that reaches every finite
history with positive probability, queries the original kernel there,
and refunds all rewards from that experiment.  It has value zero but
makes the combined transformation injective on full kernels.

Finally, an exhibited pair of equal-length programs does not establish
equality of \emph{shortest} description lengths.  A shorter universal
program might describe one partner.  The machine below handles every
complete description, including those entering through its unrestricted
universal branch.  It also excludes descriptions so short that part of
their apparent instruction tag is supplied by the encoded input history.

\subsection*{The proof in three steps}

First we construct an effective injection $F$ on all computable full
kernels.  It commutes with reward duality, preserves and reflects the
well-behaved class, and satisfies
\[
  V_{F\mu}^{\pi}=-V_{\mu}^{\pi}
\]
whenever either value exists.  Membership in each finite iterate
$F^j[\V]$ is characterized by uniformly computable finite tests.  The
intersection of all these images is one computable kernel $H$, which
has value zero for every policy.

Second, every other kernel has a finite number of inverse steps under
$F$.  Pair even depths with the next odd depths.  This defines a
set-theoretic involution $T$ reversing every policy's value.  There is no
need to decide the parity.  Two ordinary partial programs are arranged
so that the forward program is total precisely at even depths and the
inverse program is total precisely at odd depths.  The programs enforce
their promises by waiting for finitely many witnesses before answering
each requested history.

Third, six equal-length instruction tags provide the identity, reward
duality, and the two partial programs, each with either reward-duality
flag.  An input guard preserves prefix-freeness and ensures that a
complete transformed description has a complete original kernel.
Changing tags gives equal shortest lengths under both $D$ and $T$.
Absolute convergence then permits reindexing the original intelligence
sum by $T$, proving that it equals its own negative.

Sections~\ref{sec:framework}--\ref{sec:anti} set up the definitions and
the convergence argument.  Sections~\ref{sec:embedding} and
\ref{sec:core} establish the full-kernel injection and its inverse tests.
Sections~\ref{sec:parity}--\ref{sec:cancellation} construct the universal
machine and prove cancellation.

\section{Policies, kernels, and shortest descriptions}
\label{sec:framework}

Fix finite nonempty sets $A$ and $O$, and a finite nonempty set
$R\subseteq\Q\cap[-1,1]$ such that $r\in R$ implies $-r\in R$.
Symbols may be given disjoint type labels.  Write $E=O\times R$ for the
percept alphabet, $M=|E|$, and
\[
  \Hist=(EA)^*.
\]
An element
$h=(e_0,a_0,\ldots,e_{t-1},a_{t-1})\in\Hist$ has depth $|h|_E=t$.
The empty history is denoted by $\varnothing$.  If $e_i=(o_i,r_i)$,
the next action is chosen after $e_i$ is observed.  All time indices
start at zero.

\begin{definition}[Policies and full kernels]
\label{def:kernels}
A policy is a function $\pi:\Hist E\to\Dist(A)$, and an environment
kernel is a function $\mu:\Hist\to\Dist(E)$, where $\Dist(X)$ is the
set of rational probability vectors on the finite set $X$.  These
functions are defined at every history in their domains.  Let $\V$ be
the set of all total computable environment kernels.  Policies need not
be computable unless stated otherwise.
\end{definition}

The word ``full'' emphasizes that equality of kernels is equality at
every history.  An output at a history of probability zero remains part
of the function and of its description.  This convention agrees with
the source definitions in \cite{AlexanderHutter2021}.

The usual alternating sampling procedure gives finite interaction laws
for any $\mu$ and $\pi$.  Let
\begin{equation}
  Y_\mu^\pi(T)=\Exp_{\mu,\pi}\!\left[\sum_{n=0}^{T-1}r_n\right],
  \qquad T\in\N,
  \label{eq:finite-value}
\end{equation}
with $Y_\mu^\pi(0)=0$.  Every such expectation is a finite sum.  When
the limit exists as a real number, put
$V_\mu^\pi=\lim_{T\to\infty}Y_\mu^\pi(T)$.

\begin{definition}[Well-behaved environments]
\label{def:well-behaved}
Let
\[
  \W=\left\{\mu\in\V:
       \text{for every policy }\pi,
       \ V_\mu^\pi\text{ exists and belongs to }[-1,1]\right\}.
\]
The quantifier includes all rational-output policies, whether computable
or not.  No condition is imposed on the convergence of the realized
infinite reward series or on the expectations of its absolute values.
\end{definition}

For example, independent rewards from a symmetric distribution have
zero finite expected totals, even when their realized partial sums do
not converge.  Such a tail is permitted.  It will let the construction
continue forever without needing $0\in R$.

\begin{definition}[Reward duality]
\label{def:duality}
For a history or percept string $h$, let $Dh$ replace each reward $r$
by $-r$ and leave observations and actions unchanged.  Define
\[
  (D\mu)(o,r\mid h)=\mu(o,-r\mid Dh),
  \qquad
  (D\pi)(a\mid he)=\pi(a\mid D(he)).
\]
\end{definition}

\begin{lemma}[Duality]
\label{lem:duality}
The operation $D$ is an effective involution on $\V$, preserves $\W$,
and satisfies
\begin{equation}
  Y_{D\mu}^{\pi}(T)=-Y_\mu^{D\pi}(T)
  \label{eq:duality-values}
\end{equation}
for every finite horizon.  Consequently the two corresponding limiting
values exist simultaneously and are negatives.
\end{lemma}

\begin{proof}
Induct on the length of a finite interaction string.  Its probability
under $(D\mu,\pi)$ equals the probability of its reward-reversed string
under $(\mu,D\pi)$: the environment factors agree by definition, and
the action factors agree by the definition of $D\pi$.  Reward sums on
the paired strings have opposite signs.  This proves
\eqref{eq:duality-values}.  Applying $D$ twice restores each argument
and reward.  Computability is preserved, and $\pi\mapsto D\pi$ is a
bijection of all policies, so the definition of $\W$ is preserved.
\end{proof}

We use the precise program convention of
\cite[Definitions~8--10]{AlexanderHutter2021}.  Binary strings are
concatenated without an additional delimiter.

\begin{definition}[Universal machine]
\label{def:universal}
A prefix-free machine is a partial computable function
$U:\{0,1\}^*\rightharpoonup\{0,1\}^*$ whose halting domain is
prefix-free.  It is universal if, for every prefix-free partial
computable function $f$, there is a binary string $p$ such that
$U(px)=f(x)$ for every binary string $x$, including agreement on
undefinedness.
\end{definition}

\begin{definition}[Encoding, complexity, and symmetry]
\label{def:complexity}
An RL encoding is a computable injective map
\[
  \enc:\Hist\sqcup\Dist(E)\longrightarrow\{0,1\}^*
\]
whose range is prefix-free.  It is suffix-free if no distinct code is a
suffix of another code.  Fix such an encoding.  For $\mu\in\V$, let
\begin{equation}
  K_U(\mu)=\min\left\{|p|:
     U(p\enc(h))=\enc(\mu(\cdot\mid h))
     \text{ for every }h\in\Hist\right\}.
  \label{eq:complexity}
\end{equation}
A universal machine is symmetric for this encoding if
$K_U(D\mu)=K_U(\mu)$ for every $\mu\in\V$.
\end{definition}

The minimum exists for every computable kernel.  Indeed, the function
that maps $\enc(h)$ to $\enc(\mu(\cdot\mid h))$ and diverges on all
other strings is partial computable with prefix-free domain.  To
recognize an encoded history it suffices to enumerate histories until
its code is found.  Universality supplies a program for this function.
No decidability assumption about the encoding's range is needed.

\begin{definition}[Universal intelligence and its sign]
\label{def:intelligence}
For a fixed encoding and universal machine, the intelligence of a policy
is
\begin{equation}
  \Upsilon_U(\pi)=
       \sum_{\mu\in\W}2^{-K_U(\mu)}V_\mu^\pi.
  \label{eq:intelligence}
\end{equation}
Its strong sign is $\sgn(\Upsilon_U(\pi))\in\{-1,0,1\}$, with
$\sgn(0)=0$.  A weak sign procedure is required to halt and give the
correct sign when the score is nonzero; its behavior at zero is unrestricted.
\end{definition}

The summation in \eqref{eq:intelligence} is over distinct full kernels,
not over programs and not over behavioral equivalence classes.  We
will give an explicit absolute-convergence bound for this exact program
convention in Lemma~\ref{lem:kraft}.

\begin{theorem}[Zero universal intelligence]
\label{thm:main}
For every choice of finite alphabets above and every fixed suffix-free
RL encoding, there exists an ordinary prefix-free universal Turing
machine $U$ satisfying
\[
  K_U(D\mu)=K_U(\mu)\quad(\mu\in\V),
  \qquad
  \Upsilon_U(\pi)=0\quad\text{for every policy }\pi.
\]
Consequently its strong sign, and hence its weak sign, is computable
on programs promised to define total rational-output policies: the
algorithm always returns zero.
\end{theorem}

The proof occupies the following sections.  All algorithms used in the
construction are ordinary partial computations.  In particular, the
machine has no oracle for totality, for membership in $\W$, or for
equality of environment kernels.

\section{A cutoff-and-refund transformation}
\label{sec:anti}

We first construct an environment $Q\mu$ whose limiting value, when
defined, is $-2V_\mu^\pi$ for each fixed policy $\pi$.  The factor two
will be removed by a fair mixture with a zero-valued environment.

Let $u$ be the uniform law on $E$, and let $u_R$ be the uniform law on
$R$.  Both have mean reward zero.  Fix $o_*\in O$ and put
$\tau=\delta_{o_*}\otimes u_R$.  The laws $u$ and $\tau$ are invariant
under reward reversal.  Set $\eps=1/4$, and independently choose a
cutoff $N$ with
\begin{equation}
  p_0=\Prob(N=0)=\frac34,
  \qquad p_n=\Prob(N=n)=2^{-n-2}\quad(n\ge1).
  \label{eq:cutoff}
\end{equation}
Thus
\begin{equation}
  \Prob(N>t)=2^{-t-2},\qquad \Exp N=\frac12.
  \label{eq:cutoff-tail}
\end{equation}
For $n\ge1$ let $\ell_n=2^{n+2}$.  At $n=0$, independently choose
$\ell_0=1$ with probability $2/3$ and $\ell_0=2$ with probability
$1/3$.  These choices ensure
\begin{equation}
  p_n\ell_n=1\quad(n\ge1),
  \qquad p_0\Exp\ell_0=1.
  \label{eq:importance-normalization}
\end{equation}

\subsection{The physical experiment}

Conditional on $N=n$, $Q\mu$ follows this schedule.
\begin{enumerate}[label=\arabic*.,leftmargin=*]
\item At times $0,\ldots,n-1$, it produces percepts according to $\mu$
      on the actual history.  The policy's actual actions are supplied
      to $\mu$.
\item It privately samples the next source percept at that history.
      Denote its reward by $r_n$.  The percept is not shown to the
      policy, and its observation is discarded.
\item At times $n,\ldots,2n-1$, it emits observation $o_*$ and the
      rewards $-r_0,\ldots,-r_{n-1}$ in order.  This refunds the visible
      source prefix.  The interval is empty when $n=0$.
\item At times $2n,\ldots,2n+2\ell_n-1$, it emits $o_*$ and reward
      $-r_n$.
\item Thereafter it emits independent percepts with law $\tau$.
\end{enumerate}
Actions from time $n$ onward are ignored by this branch.  In particular,
the private source query occurs before an action at time $n$ could be
needed.  Up to that query the policy has seen its original source
history, so the sampled $r_n$ has its original distribution under
$(\mu,\pi)$.

Every reward emitted by the experiment belongs to $R$.  The construction
uses sign reversal and repetition, rather than a larger reward alphabet.
The payment length may have infinite expectation, which is permitted.
Each finite horizon still contains only finitely many bounded rewards.

\begin{example}[One cutoff branch]
If $N=1$, the experiment shows the first source reward $r_0$, privately
samples $r_1$, refunds $r_0$ at time $1$, and then pays $-r_1$ for
$2\ell_1=16$ steps.  The probability of this cutoff is $p_1=1/8$.
After the refund its contribution to the limiting expected total is
therefore $-2\Exp r_1$.  Every cutoff index has the same expected
coefficient $-2$.  The proof below justifies their combined limiting
effect when the original expected series is merely convergent.
\end{example}

\subsection{An ordinary rational kernel}

A latent experiment must be converted into a total rational-output
kernel before it can enter $\V$.  The following finite filtering
description also fixes its values at null histories.

\begin{lemma}[Finite filtering for $Q$]
\label{lem:q-computable}
There is a uniform ordinary index transformation $\mu\mapsto Q\mu$
which is extensional and total on $\V$.  Its finite interaction laws
are those of the preceding experiment for every policy.
\end{lemma}

\begin{proof}
At an input history $h$ of depth $t$, enumerate the alternatives
$N=0,\ldots,t$ and the single alternative $N>t$, whose probability is
given by \eqref{eq:cutoff-tail}.  For a cut branch enumerate the finitely
many possible privately sampled percepts, and for $N=0$ the two possible
values of $\ell_0$.

If $N>t$, all the observed percepts come from $\mu$, and the next
percept also has law $\mu(\cdot\mid h)$.  If $N=t>0$, the next
percept is the first refund.  If $N=t=0$, it is the first importance
payment.  If $N<t$, the branch is in a refund, payment, or terminal
phase determined by $N$, $\ell_N$, and the physical time.

For each alternative compute the likelihood of the percepts in $h$
when its listed actions are supplied externally.  The source part is
a product of outputs of $\mu$ at literal initial segments of $h$.
The refund and payment parts are checked against their prescribed
values.  The terminal part is a finite product of $\tau$ probabilities.
Hidden percept probabilities are included and summed over.  These
calculations use only finitely many rational operations and source
queries.  Normalize the likelihoods and take the corresponding mixture
of next-percept laws.  If the total likelihood is zero, return $\tau$.

On a total computable source all queries terminate and the output is a
legal rational probability vector.  The computation depends only on the
source values, so it is extensional.  It does not inspect the syntax or
running times of a source program to decide its validity.

For any fixed policy, each compatible latent alternative for $h$ has
the same action-probability factor
$\prod_{i<t}\pi(a_i\mid h_i e_i)$, where $h_i$ is the prefix before
$e_i$.  These factors cancel in a conditional next-percept law whenever
the history has positive probability.  Equivalently, induction on
finite histories verifies that the returned kernel has the latent
experiment's finite laws.  The stipulated values at zero-likelihood
histories cannot change those laws.  This induction does not require a
computable policy.
\end{proof}

\subsection{Exact finite-horizon coefficients}

Fix an arbitrary policy and write
\begin{equation}
  b_n=\Exp_{\mu,\pi}r_n,\qquad |b_n|\le1,
  \qquad S_m=\sum_{n=0}^m b_n.
  \label{eq:source-sequence}
\end{equation}
Let $B_T$ be the copied-prefix balance, including its refunds, after
the first $T$ physical rewards in the experiment.  Conditional on
$N=n$, it is zero once $T\ge2n$, and $|B_T|\le n$ at every horizon.
Since $\Exp N=1/2$, dominated convergence gives
\begin{equation}
  \Exp B_T\longrightarrow0.
  \label{eq:refund-vanishes}
\end{equation}
Every terminal reward has conditional mean zero, so the terminal part
has expected contribution zero at each finite horizon.

The importance payments therefore give the exact identity
\begin{equation}
  Y_{Q\mu}^{\pi}(T)=\Exp B_T-\sum_{n\ge0}a_{T,n}b_n,
  \label{eq:q-decomposition}
\end{equation}
where
\begin{align}
  a_{T,n}&=\min\{2,\,2^{-n-2}(T-2n)_+\}, &&n\ge1,
     \label{eq:q-coefficients}\\
  a_{T,0}&=p_0\Exp\min\{2\ell_0,T\}. &&
     \label{eq:q-coefficient-zero}
\end{align}
Here $x_+=\max\{x,0\}$.  Only finitely many coefficients are nonzero
at a finite horizon.  The factorization by $b_n$ is legitimate because
the source prefix and the private sample have their original joint
law, and the cutoff and copy counts were chosen independently.

\begin{proposition}[Convergence reflection]
\label{prop:q-convergence}
For every $\mu\in\V$ and every policy $\pi$, the sequence
$Y_{Q\mu}^{\pi}(T)$ converges if and only if $Y_\mu^\pi(T)$ converges.
When they converge,
\[
  V_{Q\mu}^{\pi}=-2V_\mu^\pi.
\]
\end{proposition}

\begin{proof}
For $T\ge4$ put $q_{T,n}=a_{T,n}/2$.  As a function of $n$, these
numbers are nonnegative, nonincreasing, eventually zero, and have
$q_{T,0}=1$.  For $n\ge1$, both $2^{-n-2}$ and $(T-2n)_+$ are
nonnegative nonincreasing functions of $n$; their product, followed by
truncation, has the same property.  For every fixed $n$,
$q_{T,n}\to1$ as $T\to\infty$.

Finite summation by parts gives
\begin{equation}
  \sum_{n\ge0}q_{T,n}b_n
   =\sum_{n\ge0}(q_{T,n}-q_{T,n+1})S_n.
  \label{eq:summation-parts}
\end{equation}
The coefficients on the right are nonnegative, sum to one, and tend
to zero at each fixed index.  If $S_n\to s$, choose a fixed index
beyond which $|S_n-s|<\eta$.  The tail error in
\eqref{eq:summation-parts} is at most $\eta$, while the finite initial
error tends to zero.  Thus the sum tends to $s$.  Equations
\eqref{eq:refund-vanishes} and \eqref{eq:q-decomposition} imply
$Y_{Q\mu}^{\pi}(T)\to-2s$.

For the converse, use the particular horizons
\[
  T_m=2^{m+3}\qquad(m\ge4),
\]
and define
\begin{equation}
  G_m=S_m+\sum_{n>m}2^{m-n}b_n.
  \label{eq:geometric-average}
\end{equation}
The tail converges absolutely because $|b_n|\le1$.  Let
$w_{m,n}=1$ for $n\le m$ and $w_{m,n}=2^{m-n}$ for $n>m$.
For $1\le n<m$, the coefficient $a_{T_m,n}$ is saturated at two.
At the largest such index this follows from
$2^{m+3}-2(m-1)\ge2^{m+2}$; the smaller indices only increase the
untruncated coefficient.  The zeroth coefficient is also two.

For $n\ge m$, comparison with
$2^{-n-2}T_m=2^{m+1-n}=2w_{m,n}$ gives
\[
  |a_{T_m,n}-2w_{m,n}|\le 2n\,2^{-n-2}.
\]
When $T_m\le2n$ the actual coefficient is zero, and the same bound
holds because $2^{-n-2}T_m\le2n\,2^{-n-2}$.  Therefore
\begin{equation}
  \sum_{n\ge0}|a_{T_m,n}-2w_{m,n}|
      \le(m+1)2^{-m}\longrightarrow0.
  \label{eq:coefficient-error}
\end{equation}
It follows from \eqref{eq:q-decomposition} that
\begin{equation}
  Y_{Q\mu}^{\pi}(T_m)+2G_m\longrightarrow0.
  \label{eq:q-geometric-subsequence}
\end{equation}

The geometric averages have the exact inverse relation
\begin{equation}
  S_m=2G_m-G_{m+1}.
  \label{eq:geometric-inverse}
\end{equation}
For $b_n$ with $n\le m$ the coefficient on the right is $2-1=1$;
for $b_{m+1}$ it is $2(1/2)-1=0$; and all subsequent geometric
coefficients cancel.  If the full sequence
$Y_{Q\mu}^{\pi}(T)$ tends to $y$, equation
\eqref{eq:q-geometric-subsequence} gives $G_m\to-y/2$, and
\eqref{eq:geometric-inverse} gives $S_m\to-y/2$.  Since
$Y_\mu^\pi(T)=S_{T-1}$ for $T\ge1$, this proves the converse and
the asserted relation between limits.
\end{proof}

The reverse implication uses the specific identity
\eqref{eq:geometric-inverse}.  A general regular summation method need
not reflect convergence.  No absolute convergence of
$\sum_n b_n$, no pathwise reward convergence, and no bound on the
intermediate $S_m$ were assumed in the reverse argument.

\section{Preserving every entry of the source kernel}
\label{sec:embedding}

The transformation $Q$ reads its source only along a copied source
trajectory.  A source entry at an unreachable history need not affect
its output.  We now add a zero-valued experiment that exposes those
entries.  A fair choice between the two experiments will also change
the factor $-2$ into $-1$.

\subsection{A zero-valued probe}

Using the same cutoff law \eqref{eq:cutoff}, define $P\mu$ by the
following experiment, conditional on $N=n$.
\begin{enumerate}[label=\arabic*.,leftmargin=*]
\item At times $0,\ldots,n-1$, emit independent percepts with law $u$.
      Accept the policy's actual actions.
\item At time $n$, query $\mu$ at the actual observed history and emit
      the resulting percept, including its observation.
\item At times $n+1,\ldots,2n+1$, emit $o_*$ and the negatives of all
      $n+1$ rewards already emitted, in their original order.
\item Thereafter emit independent percepts with law $u$.
\end{enumerate}
Actions after the probe are ignored.  The source query is made at a
history generated by the uniform experiment, which can be a null
history for $\mu$.

\begin{lemma}[The probe]
\label{lem:probe}
The map $\mu\mapsto P\mu$ is an effective extensional transformation
from $\V$ to $\V$.  For every $\mu\in\V$ and every policy,
\begin{equation}
  \lim_{T\to\infty}Y_{P\mu}^{\pi}(T)=0.
  \label{eq:probe-zero}
\end{equation}
Moreover, the likelihood of any history $h$ of depth $t$, with its
actions externally supplied, is at least
\begin{equation}
  \Prob(N\ge t)M^{-t}>0.
  \label{eq:probe-support}
\end{equation}
\end{lemma}

\begin{proof}
The finite filter enumerates $N=0,\ldots,t$ and $N>t$.  A branch
before its probe has uniform percept likelihood.  At the probe its
source query is at the literal prefix of the input history.  Its
refund and terminal likelihoods are finite rational products.  Thus
the same finite calculation as in Lemma~\ref{lem:q-computable}
produces a rational kernel, uniformly and extensionally on $\V$.
The policy factors cancel, so the calculation realizes the experiment
for every policy.

Conditional on $N=n$, the sum of rewards through the probe is
completely refunded by time $2n+1$.  Before that refund is complete,
the remaining balance has absolute value at most $n+1$.  It is zero
at every subsequent horizon apart from the terminal tail.  Since
$\Exp(N+1)=3/2$, dominated convergence makes the expected pending
balance tend to zero.  The terminal tail has expected total zero at
each finite horizon.  This proves \eqref{eq:probe-zero}, without
requiring any convergence assumption on the source.

On the event $N\ge t$, all the first $t$ percepts are independent
uniform samples.  Their likelihood is $M^{-t}$ for any supplied
action string.  This gives \eqref{eq:probe-support}.
\end{proof}

\subsection{The combined transformation}

For a kernel $\rho$ and a history
$h=(e_0,a_0,\ldots,e_{t-1},a_{t-1})$, write
\begin{equation}
  L_\rho(h)=\prod_{i=0}^{t-1}\rho(e_i\mid h_i),
  \qquad L_\rho(\varnothing)=1,
  \label{eq:likelihood}
\end{equation}
where $h_i$ is the prefix of depth $i$.  This is the percept likelihood
when the displayed actions are supplied externally.  It is distinct
from the probability of $h$ under a specified policy, which also
contains action factors.

Define $F\mu$ by a fair hidden coin, chosen once at the start: run
$Q\mu$ on one outcome and $P\mu$ on the other.  Its conditional kernel
is
\begin{equation}
 (F\mu)(\cdot\mid h)=
 \frac{L_{Q\mu}(h)(Q\mu)(\cdot\mid h)
       +L_{P\mu}(h)(P\mu)(\cdot\mid h)}
      {L_{Q\mu}(h)+L_{P\mu}(h)}.
 \label{eq:f-mixture}
\end{equation}
The denominator is positive at every history by
\eqref{eq:probe-support}.  This is a mixture of whole experiments.
Taking a fresh pointwise average of their conditional outputs at each
history would in general describe a different process.

\begin{proposition}[Value, support, and duality of $F$]
\label{prop:f-basic}
The transformation $F:\V\to\V$ is effective and extensional.  Every
$F\mu$ has full support at every history.  For every $\mu\in\V$ and
every policy $\pi$,
\begin{equation}
 \begin{split}
  V_{F\mu}^{\pi}\text{ exists}
      &\quad\Longleftrightarrow\quad V_\mu^\pi\text{ exists},\\
  V_{F\mu}^{\pi}&=-V_\mu^\pi
       \quad\text{when these limits exist}.
 \end{split}
 \label{eq:f-antivalue}
\end{equation}
Consequently
\begin{equation}
  F\mu\in\W\quad\Longleftrightarrow\quad\mu\in\W.
  \label{eq:f-w-reflection}
\end{equation}
On full kernels, $FD=DF$.
\end{proposition}

\begin{proof}
Formula \eqref{eq:f-mixture} is an ordinary finite rational computation
on every valid source.  At a history of depth $t$, the $P$ branch has
a positive contribution from the event $N>t$, on which the next
percept is uniform.  Every output coordinate is therefore positive.

Conditioning on the initial fair coin gives an identity of finite
interaction laws, and hence
\begin{equation}
  Y_{F\mu}^{\pi}(T)=
      \frac12Y_{Q\mu}^{\pi}(T)+\frac12Y_{P\mu}^{\pi}(T)
  \label{eq:f-finite-mixture}
\end{equation}
for every $T$.  Lemma~\ref{lem:probe} and
Proposition~\ref{prop:q-convergence} prove
\eqref{eq:f-antivalue} in both directions.  Quantify over all policies
and use the symmetry of $[-1,1]$ to obtain
\eqref{eq:f-w-reflection}.  In particular, a source limit outside this
interval remains outside it after negation.  The auxiliary $Q\mu$
need not itself belong to $\W$.

For duality, match latent experiments after reversing every reward in
the input and output.  Cutoffs, copy counts, and the component coin
have unchanged laws.  In $Q$, the copied rewards, refunds, and
importance payments all reverse sign.  In $P$, the uniform prefix is
invariant, the source query is made at the reversed history, its
sampled percept is reversed, and every refund is reversed.  Both
terminal laws are invariant.  Thus the unnormalized likelihoods and
next-percept vectors transform equivariantly.  The null convention
for $Q$ is also invariant.  Normalizing in
\eqref{eq:f-mixture} proves $F(D\mu)=D(F\mu)$ at every history,
including histories null for the source.
\end{proof}

\subsection{Recovering the initial distribution}

Write $q=\mu(\cdot\mid\varnothing)$, and define the stochastic
linear map
\begin{equation}
  (Jq)(o,r)=\one_{\{o=o_*\}}\sum_{v\in O}q(v,-r).
  \label{eq:initial-reversal}
\end{equation}
It discards the observation and reverses the reward.  At the initial
step, the $Q$ branch has law $\tfrac14q+\tfrac34Jq$, and the $P$
branch has law $\tfrac34q+\tfrac14u$.  Therefore the initial-row map
of $F$ is
\begin{equation}
  F_0(q)=\frac12q+\frac38Jq+\frac18u.
  \label{eq:initial-map}
\end{equation}

\begin{lemma}[Initial inverse and contraction]
\label{lem:initial-inverse}
The map $F_0$ is injective and has an exact rational inverse on its
image.  In $\ell^1$ distance it is a contraction with factor at most
$7/8$.
\end{lemma}

\begin{proof}
If $v=F_0(q)$, then for $o\ne o_*$,
\[
  q(o,r)=2v(o,r)-\frac14u(o,r).
\]
These entries determine
$s(r)=\sum_{o\ne o_*}q(o,r)$.  Put
\[
  d(r)=v(o_*,r)-\frac38s(-r)-\frac18u(o_*,r).
\]
For each pair $r,-r$ with $r\ne0$, the remaining entries solve
\[
 \begin{pmatrix}d(r)\\d(-r)\end{pmatrix}
   =\begin{pmatrix}1/2&3/8\\3/8&1/2\end{pmatrix}
     \begin{pmatrix}q(o_*,r)\\q(o_*,-r)\end{pmatrix}.
\]
The determinant is $7/64$.  If $0\in R$, its remaining equation is
$d(0)=(7/8)q(o_*,0)$.  Thus the inverse is unique and rational.
An arbitrary rational target can be checked for image membership at
this row by recovering the vector and checking nonnegativity and
total mass one.

A stochastic linear map does not increase $\ell^1$ distance:
$\norm{Jq-Jq'}_1\le\norm{q-q'}_1$.  Applying
\eqref{eq:initial-map} gives
$\norm{F_0(q)-F_0(q')}_1\le(7/8)\norm{q-q'}_1$.
\end{proof}

\subsection{A positive triangular coefficient at every history}

At a nonempty history $h$ of depth $t$, all source values at proper
prefixes of $h$ precede the unknown vector $\mu(\cdot\mid h)$.
The unnormalized next-percept law of $F\mu$ depends on that vector
with a strictly positive scalar coefficient.  This is the point of
the uniform probe.

\begin{lemma}[Triangular form]
\label{lem:triangular}
For each $h$ of depth $t\ge1$, there are a scalar $Z_h>0$, a scalar
$\gamma_h>0$, and a nonnegative vector $C_h$, all determined by the
proper source-prefix values, such that
\begin{equation}
  Z_h(F\mu)(\cdot\mid h)=
       \gamma_h\mu(\cdot\mid h)+C_h.
  \label{eq:triangular}
\end{equation}
Specifically,
\begin{align}
  Z_h&=L_{F\mu}(h),
     \label{eq:z-definition}\\
  \gamma_h&=\frac12\,2^{-t-2}L_\mu(h)
             +\frac12p_tM^{-t},
     \label{eq:gamma}\\
  C_h(e)&\ge\delta_t u(e),
  \qquad\delta_t=2^{-t-3}M^{-t}>0.
     \label{eq:delta}
\end{align}
Moreover,
\begin{equation}
  Z_h\ge\frac12\Prob(N\ge t)M^{-t},\qquad
  0<\frac{\gamma_h}{Z_h}\le1-\delta_t<1.
  \label{eq:triangular-bounds}
\end{equation}
The coefficients extend continuously to arbitrary real probability
vectors at the proper prefixes.
\end{lemma}

\begin{proof}
Consider the joint likelihood of the observed prefix and the next
percept, with the action string externally supplied.  In the $Q$
branch, the event $N>t$ has weight $2^{-t-2}$, copies the source
prefix with likelihood $L_\mu(h)$, and next samples
$\mu(\cdot\mid h)$.  It contributes the first term in
\eqref{eq:gamma}, including the fair-mixture factor $1/2$.

Every other $Q$ contribution depends only on proper source prefixes.
For $N<t$ the private query was at a proper prefix.  For $N=t>0$,
the next output is the first refund, which does not depend on the
current private sample.  That sample's probabilities sum to one,
so the current source vector disappears from this contribution.
This is why the initial row was treated separately.

In the $P$ branch, the event $N=t$ generates the prefix uniformly and
then samples $\mu(\cdot\mid h)$.  Its contribution is the second
term in \eqref{eq:gamma}.  If $N<t$, the probe query was at a proper
prefix.  If $N>t$, both the prefix and the next percept are uniform.
The last event contributes $\delta_tu$ and is included in $C_h$.
This proves \eqref{eq:triangular}--\eqref{eq:delta}.

The prefix mass $Z_h$ depends only on proper source prefixes.  It is
at least half the lower bound in \eqref{eq:probe-support}, and it is
at most one.  Summing \eqref{eq:triangular} over the next percept
gives $Z_h=\gamma_h+\sum_eC_h(e)$.  The second term in
\eqref{eq:gamma} is positive even if $L_\mu(h)=0$.  Since
$\sum_eC_h(e)\ge\delta_t$ and $Z_h\le1$, the remaining bounds
follow.

Finally, $Z_h$, $\gamma_h$, and $C_h$ can be computed directly as
finite unnormalized latent joint-law polynomials in the proper
source-prefix entries.  The current private sample in the $N=t$
refund branch is summed out before this computation.  There is no
division by a possibly zero $Q$-likelihood.  The only division in the
conditional formula is by $Z_h$, whose lower bound is uniform over
all proper-prefix probability vectors.  Thus the normalized
coefficients are continuous on the finite product of real
probability simplices, including its boundary.
\end{proof}

In particular, $F$ has retained a source entry even when
$L_\mu(h)=0$.  After its proper prefixes have been recovered, the
entry is obtained from a target $\nu=F\mu$ by
\begin{equation}
  \mu(\cdot\mid h)=
       \frac{Z_h\nu(\cdot\mid h)-C_h}{\gamma_h}.
  \label{eq:full-inverse}
\end{equation}
Together with Lemma~\ref{lem:initial-inverse}, this proves injectivity
on full kernels.  Equality of reachable interaction laws alone is
not being used as a substitute for kernel equality.

\section{Finite inverse tests and the zero backward core}
\label{sec:core}

\subsection{Testing a finite inverse table}

\begin{proposition}[Uniform inverse tests]
\label{prop:inverse-tests}
There is an ordinary partial inverse interpreter $I$ which is total
and valid exactly on $F[\V]$, with $I(F\mu)=\mu$.  There are uniform
finite rational tests $R(j,\nu,m)$, terminating for every
$\nu\in\V$ and $j,m\in\N$, such that
\begin{equation}
  \nu\in F^j[\V]\quad\Longleftrightarrow\quad
       \forall m\ R(j,\nu,m).
  \label{eq:image-tests}
\end{equation}
Set $R(0,\nu,m)$ to true.  A failed inverse test returns false rather
than calling an undefined inverse value.
\end{proposition}

\begin{proof}
Given a valid target $\nu$ and depth $m$, recover its prospective
source table at all histories through depth $m$.  Start with the
initial inverse in Lemma~\ref{lem:initial-inverse}.  Reject if its
vector is not a legal rational distribution.  Proceed in increasing
history depth, using \eqref{eq:full-inverse}, and reject if any
recovered vector is illegal.  Before each division all proper-prefix
vectors have been validated, so the strict positive bound for
$\gamma_h$ applies.  There are finitely many histories through a
fixed depth, and the computation terminates on every valid target,
whether or not it lies in the image.

An actual image passes all these tests.  Conversely, if all depths
pass, the recovered tables agree on overlaps because each row is
uniquely determined by the target row and previously recovered
prefixes.  They define a total computable rational kernel $\mu$.
The initial equation and \eqref{eq:triangular} show $F\mu=\nu$ at
every history.  Thus all finite tests pass exactly on $F[\V]$.

To compute $I\nu$ at a requested history, perform the reconstruction
through that history's depth.  Return the recovered vector if it
passes, and diverge if it fails.  An image has a total inverse.  If a
target is not an image, some depth fails, so the interpreter is not
total.  This gives the asserted partial inverse without deciding
image membership.

For the $j$-fold test, reconstruct the finite depth-$m$ table $j$
successive times, rejecting at the first failure.  Each reconstruction
uses only a table of that same depth.  This defines a terminating
test $R(j,\nu,m)$.  If $\nu=F^j\mu$, the corresponding $j$ inverse
tables all pass.  If every depth passes, consistency gives successively
$j$ total computable kernels linked by $F$; equivalently one may
induct on $j$ using the one-step result.  The last is a preimage of
$\nu$ under $F^j$.  This proves \eqref{eq:image-tests}.
\end{proof}

This proposition gives a universal conjunction of finite tests for
image membership.  It does not give an image-membership decision
procedure.  Its strength is that each \emph{individual} test is
total on every valid target, including targets outside the image.

\subsection{The fixed kernel}

\begin{proposition}[A computable fixed point]
\label{prop:fixed-point}
There is a unique fixed kernel $H$ for $F$ on $\V$.  It is rational,
computable, and has full support.  For every finite depth $m$,
\begin{equation}
  \sup_{\mu\in\V}\ \max_{|h|_E\le m}
       \norm{(F^n\mu)(\cdot\mid h)-H(\cdot\mid h)}_1
       \longrightarrow0
       \quad(n\to\infty).
  \label{eq:uniform-convergence}
\end{equation}
\end{proposition}

\begin{proof}
The initial-row map is a contraction by
Lemma~\ref{lem:initial-inverse}.  Its unique fixed distribution is
explicitly
\begin{equation}
  H(\cdot\mid\varnothing)=
  \left(\frac34\delta_{o_*}+\frac14\operatorname{Unif}(O)\right)
       \otimes u_R.
  \label{eq:initial-fixed}
\end{equation}
Substitution into \eqref{eq:initial-map} verifies this formula.

Suppose the fixed values at all proper prefixes of a nonempty history
$h$ have been determined.  With those coefficients fixed, write
\eqref{eq:triangular} in the form
\begin{equation}
  (F\mu)(\cdot\mid h)=a_h\mu(\cdot\mid h)+c_h,
  \qquad a_h=\frac{\gamma_h}{Z_h},\quad c_h=\frac{C_h}{Z_h}.
  \label{eq:affine-row}
\end{equation}
Here $0<a_h\le1-\delta_{|h|_E}$, $c_h$ is nonnegative, and
$\sum_ec_h(e)=1-a_h$.  The unique fixed row is therefore
\begin{equation}
  H(\cdot\mid h)=\frac{c_h}{1-a_h}.
  \label{eq:fixed-row}
\end{equation}
It is a rational probability vector.  Recursion over the finitely many
prefixes needed for any requested history makes $H$ a total computable
kernel.  The positive uniform contribution in
\eqref{eq:delta} makes every coordinate positive.  The same
recursion proves uniqueness of a fixed kernel, even among all
real-valued full probability kernels.

It remains to prove the uniform convergence needed later.  Work at
each fixed depth with the finite product of \emph{real} probability
simplices; this is compact.  The coefficient maps in
\eqref{eq:affine-row} are continuous on that whole product by
Lemma~\ref{lem:triangular}, and hence uniformly continuous.  We do
not assume that the countable set $\V$ itself is compact.

At depth zero, the contraction factor $7/8$ gives uniform convergence
from the bounded initial simplex.  Inductively suppose all proper
prefix rows of $F^n\mu$ converge uniformly in $\mu$ to those of $H$.
For a depth-$t$ history $h$, let $a_h^{(n)}$ and $c_h^{(n)}$ be the
coefficients calculated from these rows.  Uniform continuity and the
induction hypothesis give
$a_h^{(n)}\to a_h^*$ and $c_h^{(n)}\to c_h^*$ uniformly, where
the starred coefficients are those of $H$.  Subtract the fixed-row
equation to obtain
\begin{align*}
 &\norm{(F^{n+1}\mu)(\cdot\mid h)-H(\cdot\mid h)}_1\\
 &\quad\le (1-\delta_t)
       \norm{(F^n\mu)(\cdot\mid h)-H(\cdot\mid h)}_1
       + e_t(n),
\end{align*}
where $e_t(n)\to0$ uniformly in $\mu$.  There are only finitely
many histories at depth $t$, so the same error sequence can be used
for all of them.  For example, it may be taken to bound
$|a_h^{(n)}-a_h^*|+\norm{c_h^{(n)}-c_h^*}_1$.

If a bounded nonnegative sequence satisfies
$x_{n+1}\le\rho x_n+e_n$ with $0\le\rho<1$ and $e_n\to0$,
then
\[
  x_n\le\rho^nx_0+\sum_{k<n}\rho^{n-1-k}e_k\longrightarrow0.
\]
To see the last limit, split the sum at a fixed index after which
$e_k$ is arbitrarily small.  The finite initial part decays
geometrically; the later part is at most
$\sup_{k\ge K}e_k/(1-\rho)$.  Applying this estimate uniformly
proves convergence at depth $t$ and completes the induction.
\end{proof}

\begin{proposition}[The backward core is one zero-valued kernel]
\label{prop:core}
The kernel $H$ is independent of the action entries of its history,
is self-dual, and belongs to $\W$.  Every policy has value zero in it.
Moreover,
\begin{equation}
  \bigcap_{j\ge0}F^j[\V]=\{H\}.
  \label{eq:core}
\end{equation}
\end{proposition}

\begin{proof}
Both latent components preserve action independence of the source:
their schedules and known percept laws ignore actions, and any source
query then depends only on the percept entries.  Their finite filters
also ignore the listed actions.  Thus $F$ preserves action independence.
It preserves self-duality because it commutes with $D$.

Start iteration at the independent uniform kernel.  Each iterate is
action-independent and self-dual.  These are closed coordinate
identities, so \eqref{eq:uniform-convergence} passes them to $H$.
For an action-independent kernel, summing over policy actions leaves
the same finite percept law for every policy.  Self-duality makes
that percept law invariant under reversing all rewards.  The expected
sum of rewards at each finite horizon is therefore zero.  Thus
$V_H^\pi=0$ for every policy and $H\in\W$.  Both properties are
used here: self-duality alone would relate a policy to its dual and
would not force zero for each fixed reward-sensitive policy.

If $\nu\in\bigcap_jF^j[\V]$, for each $j$ choose any kernel
$\mu_j$ with $F^j\mu_j=\nu$.  For each fixed depth $m$,
\eqref{eq:uniform-convergence} forces the table of $\nu$ through
that depth to equal the table of $H$.  Hence $\nu=H$ as a full
kernel.  This argument requires neither coherent preimages across
$j$ nor an effective choice of them.  Conversely, $F(H)=H$ puts
$H$ in every image.  This proves \eqref{eq:core}.
\end{proof}

The construction has now supplied both the analytic and the effective
facts needed by the universal machine.  The analytic part is the
same-policy sign reversal with exact convergence reflection.  The
effective part consists of a full-kernel injection and terminating
finite tests for each finite inverse depth.  The single remaining
kernel in the infinite backward core contributes zero.

\section{Pairing inverse depths with partial programs}
\label{sec:parity}

The injection $F$ organizes kernels into forward chains, apart from its
single fixed kernel.  Consecutive kernels in a chain have opposite
values for every policy.  Pairing consecutive depths is straightforward
as a set-theoretic operation.  The issue is how to make a universal
language respect that pairing without deciding inverse depth.

\begin{definition}[Backward depth]
\label{def:depth}
For $j\in\N$ and $\nu\in\V$, write
\[
  D_j(\nu)\quad\Longleftrightarrow\quad\nu\in F^j[\V].
\]
For $\nu\ne H$, its backward depth is
\[
  d(\nu)=\max\{j:D_j(\nu)\}.
\]
\end{definition}

The images are nested: $F^{j+1}[\V]\subseteq F^j[\V]$.
Proposition~\ref{prop:core} therefore guarantees that the displayed
maximum exists for every $\nu\ne H$.  No algorithm for computing it
is part of this definition.

\begin{lemma}[Adjacent-depth pairing]
\label{lem:pairing}
For $\nu\ne H$,
\[
  d(F\nu)=d(\nu)+1,
  \qquad
  d(I\nu)=d(\nu)-1\quad\text{if }d(\nu)>0.
\]
Define
\begin{equation}
  T\nu=
  \begin{cases}
    H,&\nu=H,\\
    F\nu,&d(\nu)\text{ is even},\\
    I\nu,&d(\nu)\text{ is odd}.
  \end{cases}
  \label{eq:pairing}
\end{equation}
Then $T$ is an involution of $\V$, commutes with $D$, maps $\W$
bijectively to itself, and satisfies
\begin{equation}
  V_{T\nu}^{\pi}=-V_\nu^\pi
  \qquad(\nu\in\W)
  \label{eq:t-antivalue}
\end{equation}
for every policy.
\end{lemma}

\begin{proof}
If $d(\nu)=d$, then $F\nu\in F^{d+1}[\V]$.  If it also belonged
to $F^{d+2}[\V]$, injectivity of $F$ would imply
$\nu\in F^{d+1}[\V]$, a contradiction.  This proves the first
depth identity; the inverse identity follows by cancellation in the
same way.  Injectivity also gives $F^{-1}\{H\}=\{H\}$.

The definition pairs depths $0$ and $1$, $2$ and $3$, and so on,
using the unique inverse on positive depths.  Thus $T^2=\id$.
Proposition~\ref{prop:f-basic} and its inverse implication preserve
$\W$ and negate the value at each finite-depth pair.  At $H$ the
value is zero by Proposition~\ref{prop:core}.

Since $D$ is a bijection and commutes with $F$,
$D(F^j[\V])=F^j[\V]$ for every $j$.  It therefore preserves finite
depth and the core.  Injectivity gives $ID=DI$ on $F[\V]$:
both sides are the unique preimage under $F$ of the same dual kernel.
Using these identities in \eqref{eq:pairing} proves $TD=DT$.
\end{proof}

\subsection{Two promises of the right logical form}

Define the predicates
\begin{align}
  \mathsf{Even}(\nu)
    &\quad\Longleftrightarrow\quad
       \forall n\,[D_{2n+1}(\nu)\Rightarrow D_{2n+2}(\nu)],
       \label{eq:even-predicate}\\
  \mathsf{Odd}(\nu)
    &\quad\Longleftrightarrow\quad
       D_1(\nu)\ \wedge\
       \forall n\,[D_{2n+2}(\nu)\Rightarrow D_{2n+3}(\nu)].
       \label{eq:odd-predicate}
\end{align}
At a finite depth these hold exactly for the indicated parity.  For
example, at depth $2k+1$ the $n=k$ implication in
\eqref{eq:even-predicate} fails; at depth $2k$ all its implications
hold.  The initial $D_1$ condition removes depth zero from the odd
predicate.  Both predicates hold at $H$.

The useful feature is their expression as universal families of
finite-witness searches.  From Proposition~\ref{prop:inverse-tests},
for any $a,b\in\N$,
\begin{equation}
 \begin{split}
  D_a(\nu)\Rightarrow D_b(\nu)
  \quad\Longleftrightarrow\quad
  \forall m\,\exists k\,
      [\neg R(a,\nu,k)\ \lor\ R(b,\nu,m)].
 \end{split}
 \label{eq:guard-normal-form}
\end{equation}
If $D_a$ fails, a finite failed test is a witness for every $m$.
If $D_a$ holds, no such failed test exists, and the right side
requires every test for $D_b$ to pass.  This proves both directions.
The bracketed predicate is decidable from finitely many exact outputs
of a valid $\nu$.

Fix a computable bijective enumeration
\[
  h_0,h_1,h_2,\ldots
\]
of all histories, and a computable enumeration of all pairs $(n,m)$.
The inverse index of each supplied finite history is computable, for
example by searching a length-lexicographic enumeration.  The index
in this enumeration is separate from the physical time or depth of
that history.

At a requested history $h_t$, the \emph{even guard} searches for
witnesses to the first $t+1$ obligations in
\eqref{eq:guard-normal-form}, with
$(a,b)=(2n+1,2n+2)$.  It returns only after all those witnesses have
been found.  The \emph{odd guard} does the same with
$(a,b)=(2n+2,2n+3)$ and also requires the first $t+1$ finite tests
$R(1,\nu,m)$ to pass.  If one of those explicit tests fails, the
request diverges.  These are stateless computations: each request
performs its own finite collection of searches.

\begin{lemma}[Totality of the guards]
\label{lem:guards}
On any $\nu\in\V$, the even guard returns at every requested
history exactly when $\mathsf{Even}(\nu)$ holds.  The corresponding
statement holds for the odd guard and $\mathsf{Odd}(\nu)$.
\end{lemma}

\begin{proof}
If the predicate holds, each obligation has a witness.  Every test
performed during a witness search terminates, so any finite collection
of searches finishes.  In the odd case all the additional $D_1$
tests also pass.

If an implication predicate fails, one of its enumerated $(n,m)$
obligations has no witness.  For every history index at or beyond
that obligation, the guard includes a search which never finishes.
If the odd predicate fails through $D_1$, some finite test fails and
likewise prevents all sufficiently late indexed requests from
returning.  Thus the relevant guarded procedure is not total.
\end{proof}

\begin{definition}[Partial forward and inverse programs]
\label{def:partial-branches}
Given an index for a partial source kernel $\nu$, define two partial
kernel programs $A_\nu$ and $B_\nu$ as follows.  At a requested
history, $A_\nu$ first completes the even-guard search and then returns the
literal value of $F\nu$ there.  The program $B_\nu$ first completes the odd-guard search,
then runs $I\nu$ there, and returns its result only if it is a legal
rational distribution.
\end{definition}

\begin{proposition}[Two partial implementations of the pairing]
\label{prop:partial-pairing}
For $\nu\in\V$, the program $A_\nu$ is a full valid kernel exactly
when $\mathsf{Even}(\nu)$ holds, and then equals $F\nu$.  The program
$B_\nu$ is a full valid kernel exactly when $\mathsf{Odd}(\nu)$
holds, and then equals $I\nu$.  At every finite-depth kernel exactly
one of these programs is full, and it computes $T\nu$.  At $H$ both
are full and return $H$.
\end{proposition}

\begin{proof}
The forward transformation is total on every valid source, so its
guard alone determines totality.  A true odd guard includes $D_1$,
which ensures that the inverse interpreter is total and returns a
valid kernel.  False guards prevent totality by
Lemma~\ref{lem:guards}.  The parity characterizations above and
$F(H)=I(H)=H$ give the remaining statements.
\end{proof}

For clarity, on a chain with root $\rho\notin F[\V]$ the full
branch and its output are:
\[
\begin{array}{c|cccccc}
 \text{input}&\rho&F\rho&F^2\rho&F^3\rho&F^4\rho&F^5\rho\\
 \hline
 \text{full branch}&A&B&A&B&A&B\\
 \text{output}&F\rho&\rho&F^3\rho&F^2\rho&F^5\rho&F^4\rho
\end{array}
\]
The table is a semantic statement about which programs are total.
It does not supply a decision procedure for choosing that branch.

In particular, racing $A_\nu$ and $B_\nu$ on one requested history
does not compute the involution $T$.  A false guard can pass its first
several obligations, so the wrong branch may return on some histories.
Its failure is a failure of totality over all histories.  The
universal machine will contain both partial programs.  The complexity
argument needs the existence of a correct complete description, not
an effective choice between them.

\section{An ordinary prefix-free universal machine}
\label{sec:machine}

Fix the suffix-free RL encoding from Theorem~\ref{thm:main}, and let
$U_0$ be any ordinary prefix-free universal machine.  We construct a
new machine $U$ with six instruction tags.  The identity tag leaves
$U_0$ unrestricted on arbitrary binary inputs.  The other tags
interpret an encoded environment request and apply duality or one of
the guarded programs.

\subsection{Tags that cannot consume part of an encoded input}

Let
\[
  e_* = \enc(\varnothing),\qquad L=|e_*|.
\]
The code $e_*$ is nonempty: an empty word could not belong to the
prefix-free range of this infinite injective encoding.  Let $b$ be
its first bit and $c=1-b$.  Choose six distinct three-bit words
$a_g$, indexed by
\[
  g\in\{\id,D,A,DA,B,DB\},
\]
and define
\begin{equation}
  \tagword_g=a_gc^L,\qquad m=L+3.
  \label{eq:tags}
\end{equation}
All tags have length $m$.  The machine diverges on inputs of length
less than $m$ and on inputs whose first $m$ bits are not a tag.

\begin{lemma}[Full descriptions contain a whole tag]
\label{lem:framing}
For any machine using the tag convention just specified, no program
of length less than $m$ describes a full environment.  Every full
environment description has the syntactic form $\tagword_g p$.
\end{lemma}

\begin{proof}
Suppose $q$ has length $r<m$, and evaluate it at the empty-history
code.  Put $j=m-r$.  If $j>L$, the string $qe_*$ has length less
than $m$ and the machine diverges.  If $1\le j\le L$, the final
$j$ bits of the first $m$ bits of $qe_*$ are the first $j$ bits of
$e_*$, which begin with $b$.  The final $j$ bits of every valid
tag are $c^j$, which begin with $c\ne b$.  Hence the first $m$
bits cannot form a tag, and the machine again diverges.

A full description must return a value at the empty history, so it
cannot have length less than $m$.  Its first $m$ bits must then be
a tag, which gives the claimed form.
\end{proof}

Equal-length tags alone would not prove this lemma: programs in
\eqref{eq:complexity} are arbitrary strings prepended to encoded
inputs.  The tail $c^L$ in \eqref{eq:tags} prevents a purported short
program from completing its tag with bits from $e_*$.  The encoding
itself has not been changed.

\subsection{The six branches}

Define the identity branch on every binary string $x$ by
\begin{equation}
  U(\tagword_{\id}x)=U_0(x).
  \label{eq:identity-port}
\end{equation}
There is no encoding restriction on this branch.

For any other tag, enumerate histories until one is found whose code
is a suffix of $x$.  Such a history is unique if it exists: two
different suffix codes would make the shorter code a suffix of the
longer one, contradicting suffix-freeness.  Write
\[
  x=p\enc(h).
\]
If no such history is found, the computation does not return.
Only a search is required; the encoding's range need not be decidable.

Before applying any transformation, run $U_0(x)$ and require its
output to encode a legal rational probability vector.  This is the
\emph{original-input guard}.  Decoding can be implemented by
enumerating rational probability vectors until the output code is
found; an output outside that range causes divergence.  From the
remaining program string $p$, compile the partial kernel
\begin{equation}
  \nu_p(\cdot\mid h')=
      \operatorname{decode}\bigl(U_0(p\enc(h'))\bigr),
  \label{eq:base-kernel}
\end{equation}
where undefined or invalidly encoded answers are treated as
divergence.

After the original-input guard has succeeded, the tag specifies the
following output at $h$:
\begin{center}
\begin{tabular}{ll}
\toprule
Tag & Partial kernel evaluated at $h$\\
\midrule
$D$  & $D\nu_p$\\
$A$  & $A_{\nu_p}$\\
$DA$ & $D(A_{\nu_p})$\\
$B$  & $B_{\nu_p}$\\
$DB$ & $D(B_{\nu_p})$\\
\bottomrule
\end{tabular}
\end{center}
Return the encoded rational vector when this computation terminates
validly.  For dual tags, $D$ includes evaluating the underlying partial
kernel at $Dh$ and reversing the output rewards.  Thus it uses the
same full-kernel duality as Definition~\ref{def:duality}.

Every operation just described is a partial computable operation on
binary strings.  The finite inverse tests are terminating on valid
base kernels; the witness searches may wait forever.  No branch
contains a test deciding that its base kernel is total or that a
parity predicate holds.

\begin{proposition}[Prefix-freeness and universality]
\label{prop:universal-machine}
The machine $U$ is an ordinary prefix-free universal Turing machine.
\end{proposition}

\begin{proof}
For every nonidentity tag, a halting input $\tagword_gx$ must first
pass the original-input guard, so $x\in\dom(U_0)$.  Its untagged
halting domain is therefore a subset of the prefix-free set
$\dom(U_0)$.  The identity branch has exactly that halting domain.
Distinct tags have the same length and are different, so a halting
word under one tag cannot be a prefix of a halting word under
another.  The union of the six tagged domains is prefix-free.

For any prefix-free partial computable binary function $f$, choose
a $U_0$ program $p_f$ with $U_0(p_fx)=f(x)$ for every binary $x$.
Then \eqref{eq:identity-port} gives
\[
  U(\tagword_{\id}p_fx)=f(x)
\]
on all inputs, including equality of undefinedness.  Hence $U$ is
universal in the full sense of Definition~\ref{def:universal}.
\end{proof}

The original-input guard uses $U_0(x)$ at the \emph{actual} untagged
input, even when a dual branch subsequently queries a different
history.  This is what makes each branch's untagged halting domain
a subset of one fixed prefix-free domain.

\begin{proposition}[Reflection of complete kernels]
\label{prop:totality-reflection}
If $\tagword_g p$ describes a complete valid environment under $U$,
then $\nu_p\in\V$.  For a valid base kernel, the identity and dual
branches are complete.  The $A$ and $DA$ branches are complete
exactly on $\mathsf{Even}$, and the $B$ and $DB$ branches exactly
on $\mathsf{Odd}$.  Every complete branch preserves and reflects
membership in $\W$.
\end{proposition}

\begin{proof}
On the identity branch, validity at every encoded history directly
asserts that \eqref{eq:base-kernel} is total and legal.  On any
other branch, a valid output at each $h$ requires the original-input
guard $U_0(p\enc(h))$ to halt with a legal vector.  The suffix
parse at such an input is necessarily the displayed $h$, by
uniqueness.  Thus the same conclusion holds.  In particular, a
transformation cannot turn a partial base program into a complete
environment by avoiding some of its inputs or filling in missing
entries.

On a valid base kernel the original-input guard always finishes.
Proposition~\ref{prop:partial-pairing} supplies the exact totality
conditions for $A$ and $B$.  Composing with $D$ does not change
totality, since $h\mapsto Dh$ is a bijection and reward reversal
preserves legal distributions.  The identity and dual branches
are total on every valid base kernel.

Finally, duality preserves and reflects $\W$ by
Lemma~\ref{lem:duality}.  The same is true of $F$ by
\eqref{eq:f-w-reflection}, and of its inverse on its image by that
equivalence in the opposite direction.  The guarded complete
branches compute exactly these kernels, so each complete branch
preserves and reflects $\W$.
\end{proof}

\section{Shortest descriptions and cancellation}
\label{sec:cancellation}

The preceding construction gave six partial instruction branches,
not a new definition of the intelligence weight.  The weight remains
$2^{-K_U(\mu)}$ with the minimum taken over all descriptions.
We now prove the needed equalities for these actual minima.

\begin{proposition}[Symmetry on all computable kernels]
\label{prop:shortest-duality}
For every $\nu\in\V$,
\[
  K_U(D\nu)=K_U(\nu).
\]
\end{proposition}

\begin{proof}
Take any complete $U$ description of $\nu$.  By
Lemma~\ref{lem:framing} it has the form $\tagword_g p$.  Toggle
the duality flag by the tag pairs
\[
  \id\longleftrightarrow D,\qquad
  A\longleftrightarrow DA,\qquad
  B\longleftrightarrow DB.
\]
Proposition~\ref{prop:totality-reflection} ensures a valid base
kernel and the same totality of the flagged and unflagged branches.
The replacement is a complete description of $D\nu$ of the same
length.  Taking a shortest description gives
$K_U(D\nu)\le K_U(\nu)$.  Apply the same argument to $D\nu$
and use $D^2=\id$ for the reverse inequality.
\end{proof}

This is symmetry on all of $\V$.  Neither well-behavedness nor
finite backward depth was assumed.

\begin{proposition}[Equal shortest weights for the value pairing]
\label{prop:shortest-pairing}
For every $\nu\in\W\setminus\{H\}$,
\begin{equation}
  K_U(T\nu)=K_U(\nu).
  \label{eq:shortest-pairing}
\end{equation}
\end{proposition}

\begin{proof}
Take an arbitrary complete description $\tagword_g p$ of $\nu$.
The base kernel $\mu=\nu_p$ is valid by
Proposition~\ref{prop:totality-reflection}.  It belongs to $\W$
because the complete branch reflects well-behavedness.  It has finite
backward depth: otherwise $\mu=H$, and every complete branch on it
would return $H$, contrary to $\nu\ne H$.

Suppose first that $g$ is $\id$ or $D$, with duality flag
$\sigma\in\{0,1\}$.  Then $\nu=D^\sigma\mu$.
Exactly one of $A_\mu$ and $B_\mu$ is complete, and it returns
$T\mu$.  Replace the tag with that branch, retaining the same
duality flag.  The new program returns
\[
  D^\sigma T\mu=T(D^\sigma\mu)=T\nu
\]
and has the same length.

If the tag is $A$, $DA$, $B$, or $DB$, its complete output is
$D^\sigma T\mu$ for the corresponding flag $\sigma$.  Replace
it with $\id$ or $D$, respectively.  The new complete output is
\[
  D^\sigma\mu=T(D^\sigma T\mu)=T\nu,
\]
again with the same length.  Here we used $T^2=\id$ and $TD=DT$.

Thus \emph{every} complete description of $\nu$ has a same-length
complete description of $T\nu$.  This includes any short program
through the unrestricted universal branch, regardless of how it
computes its kernel.  Applying the statement to a shortest
description gives $K_U(T\nu)\le K_U(\nu)$.  Apply it again to
$T\nu$ and use involutivity to obtain equality.
\end{proof}

The replacement of an identity tag in this proof need not be
computably chosen from the source program.  A Kolmogorov-complexity
inequality asks whether a description of a certain length exists.
Both guarded programs already belong to one ordinary machine; the
proof establishes that one of them is complete.  No noncomputable
choice is used during execution of $U$.

\begin{lemma}[Absolute convergence with the exact program convention]
\label{lem:kraft}
For the fixed encoding and any prefix-free universal machine,
\begin{equation}
  \sum_{\mu\in\W}2^{-K_U(\mu)}\le2^L,
  \qquad L=|\enc(\varnothing)|.
  \label{eq:shifted-kraft}
\end{equation}
Consequently the intelligence series is absolutely convergent for
every policy.
\end{lemma}

\begin{proof}
Choose the lexicographically first shortest description $q_\mu$ for
each $\mu\in\W$.  Distinct kernels have distinct descriptions:
one program returning a full kernel at every history cannot describe
two different functions.  The strings $q_\mu e_*$ are therefore
distinct members of $\dom(U)$.  They form a prefix-free set, so
Kraft's inequality gives
\[
  \sum_{\mu\in\W}2^{-(|q_\mu|+L)}\le1.
\]
Indeed, their binary cylinder sets are disjoint, and the cylinder of a
word of length $n$ has fair-coin measure $2^{-n}$.
Multiplying by $2^L$ proves \eqref{eq:shifted-kraft}.  For every
$\mu\in\W$ and every policy, $|V_\mu^\pi|\le1$ by definition.
Comparison with \eqref{eq:shifted-kraft} proves absolute
convergence.
\end{proof}

The shift by $L$ avoids assuming that the bare program strings
$q_\mu$ themselves form a prefix code.  Prefix-freeness is known
for complete halting inputs, and appending one fixed encoded history
puts all descriptions into that domain.  The precise bound is more
than is needed; its finiteness licenses rearrangement of the original
kernel sum.

\begin{proof}[Proof of Theorem~\ref{thm:main}]
Proposition~\ref{prop:universal-machine} proves ordinary
prefix-free universality, and
Proposition~\ref{prop:shortest-duality} proves symmetry on every
computable full kernel.  Fix any policy $\pi$.  The contribution
of $H$ to \eqref{eq:intelligence} is zero by
Proposition~\ref{prop:core}.  On the remaining kernels, $T$ is a
bijection of $\W\setminus\{H\}$, preserves the shortest length
by Proposition~\ref{prop:shortest-pairing}, and negates the value
by Lemma~\ref{lem:pairing}.  Absolute convergence from
Lemma~\ref{lem:kraft} permits reindexing, so
\begin{align*}
  \Upsilon_U(\pi)
   &=\sum_{\nu\in\W\setminus\{H\}}
         2^{-K_U(\nu)}V_\nu^\pi\\
   &=\sum_{\nu\in\W\setminus\{H\}}
         2^{-K_U(T\nu)}V_{T\nu}^\pi\\
   &=-\sum_{\nu\in\W\setminus\{H\}}
         2^{-K_U(\nu)}V_\nu^\pi
    =-\Upsilon_U(\pi).
\end{align*}
Hence $\Upsilon_U(\pi)=0$.  The argument did not assume that
$\pi$ is computable.  On a program promised to describe a
computable policy, the constant-zero algorithm therefore returns
its strong sign and also satisfies the weak-sign requirement.
\end{proof}

\section{Consequences and scope}
\label{sec:scope}

The construction works with arbitrary nonempty finite action and
observation sets and every finite symmetric rational reward set.
It applies, for example, when there are two actions and the only
rewards are $-1$ and $1$.  Zero need not be an available individual
reward: symmetric terminal randomization has zero finite expected
contribution.  Thus the conclusion does not depend on using a
degenerate action or reward alphabet.

Every computable environment retains a positive finite-complexity
weight.  The universal simulation branch remains unrestricted, and
the aggregate uses one shortest-description weight per full kernel.
The cancellation occurs because the additional program branches
force the minima of two semantically paired kernels to agree.
Neither suppressing difficult environments nor replacing the kernel
sum by a sum over descriptions is involved.

The all-policy condition in $\W$ is preserved in full.  The
convergence argument fixed an arbitrary policy, and finite filtering
used only its finite action probabilities.  The same-policy value
identity therefore holds even for noncomputable policies, and
well-behavedness is reflected as well as preserved.  This reflection
is needed when an arbitrary complete universal description is
reinterpreted by another tag: a well-behaved output must come from
a well-behaved base kernel.

The result separates the sign question from the usefulness of the
score as a ranking.  The constructed score and its absolute value
are both zero for all policies.  A strengthened question could ask
for a symmetric universal machine with computable sign and with
$\Upsilon_U(\pi)\ne0$ for at least one policy.  Another could require
two computable policies to receive different scores.  The argument
here establishes neither strengthening.  It shows that an explicit
nontriviality requirement is necessary if the existence question is
intended to rule out universal cancellation.

\section*{Provenance}

The repository owner selected the source problem and directed the
research.  OpenAI Codex developed the mathematical arguments and
exposition.  The proof is presented in full for mathematical scrutiny;
independent human mathematical review, formal verification, and
historical priority are not claimed.

\begin{thebibliography}{9}

\bibitem{AlexanderHutter2021}
Samuel Allen Alexander and Marcus Hutter.
\newblock Reward-Punishment Symmetric Universal Intelligence.
\newblock In Ben Goertzel, Matthew Ikl\'e, and Alexey Potapov, editors,
\emph{Artificial General Intelligence (AGI 2021)}, volume 13154 of
\emph{Lecture Notes in Computer Science}, pp.~1--10. Springer, Cham, 2022.
\newblock \href{https://doi.org/10.1007/978-3-030-93758-4_1}
  {doi:10.1007/978-3-030-93758-4\_1}.
\newblock Preprint: \href{https://arxiv.org/abs/2110.02450v1}
  {arXiv:2110.02450v1}.

\bibitem{LeggHutter2007}
Shane Legg and Marcus Hutter.
\newblock Universal Intelligence: A Definition of Machine Intelligence.
\newblock \emph{Minds and Machines}, 17(4):391--444, 2007.
\newblock \href{https://doi.org/10.1007/s11023-007-9079-x}
  {doi:10.1007/s11023-007-9079-x}.
\newblock Preprint: \href{https://arxiv.org/abs/0712.3329}
  {arXiv:0712.3329}.

\end{thebibliography}
\end{document}
